% ============================================================
%  Observability Is Restorability
%  Restorable Distinctions and the Geometry of Perceptual Access
% ============================================================
%  Engine: LuaLaTeX
%  Font:   TeX Gyre Pagella
% ============================================================

\documentclass[12pt, oneside]{article}

% ── Fonts & encoding ────────────────────────────────────────
\usepackage{fontspec}
\setmainfont{TeX Gyre Pagella}
\usepackage{unicode-math}
\setmathfont{TeX Gyre Pagella Math}

% ── Layout ──────────────────────────────────────────────────
\usepackage[
  top=1.25in, bottom=1.25in,
  left=1.35in, right=1.35in
]{geometry}
\usepackage{setspace}
\setstretch{1.15}
\setlength{\parskip}{0.55em}
\setlength{\parindent}{1.4em}

% ── Mathematics ─────────────────────────────────────────────
\usepackage{amsmath, amsthm, amssymb}
\usepackage{mathtools}

% Theorem environments
\theoremstyle{plain}
\newtheorem{theorem}{Theorem}[section]
\newtheorem{proposition}[theorem]{Proposition}
\newtheorem{lemma}[theorem]{Lemma}
\newtheorem{corollary}[theorem]{Corollary}

\theoremstyle{definition}
\newtheorem{definition}[theorem]{Definition}
\newtheorem{example}[theorem]{Example}

\theoremstyle{remark}
\newtheorem{remark}[theorem]{Remark}

\theoremstyle{definition}
\newtheorem{axiom}[theorem]{Axiom}

% ── Notation shorthands ─────────────────────────────────────
\DeclareMathOperator{\Obs}{Obs}
\DeclareMathOperator{\Rest}{Rest}
\DeclareMathOperator{\Reach}{Reach}
\DeclareMathOperator{\spec}{spec}
\newcommand{\ObsStab}{\mathrm{Obs}_{\mathrm{stab}}}
\newcommand{\Restf}{\Rest_{f}}
\newcommand{\Reston}{\Rest_{\omega}}
\newcommand{\CR}{C_{\!R}}
\newcommand{\IntMat}{\mathbf{I}}
\newcommand{\AO}{\mathcal{A}_O}
\newcommand{\DO}{\mathcal{D}_O}
\newcommand{\RO}{\mathcal{R}_O}
\newcommand{\bA}{\partial\mathcal{A}}
\newcommand{\Xsub}{X_S}
\newcommand{\Psub}{\mathcal{P}}
\newcommand{\Dsub}{\mathcal{D}}

% ── Section & title styling ─────────────────────────────────
\usepackage{titlesec}
\titleformat{\section}{\normalfont\large\bfseries}{{\thesection}.}{0.6em}{}
\titleformat{\subsection}{\normalfont\normalsize\bfseries}{{\thesubsection}.}{0.5em}{}

% ── Hyperref (load last) ────────────────────────────────────
\usepackage[
  colorlinks=true,
  linkcolor=black,
  citecolor=black,
  urlcolor=black
]{hyperref}

% ============================================================
\begin{document}

% ── Title block ─────────────────────────────────────────────
\begin{center}
  {\LARGE\bfseries Observability Is Restorability}\\[0.9em]
  {\large Restorable Distinctions and the Geometry of Perceptual Access}\\[1.5em]
  {\normalsize Flyxion}\\[0.3em]
  {\small\itshape Independent Researcher}\\[1.6em]
  {\small\today}
\end{center}

\vspace{0.5em}

% ── Abstract ────────────────────────────────────────────────
\begin{abstract}
\noindent
Standard theories of observation treat the observer as a receiver:
a system that acquires, records, or represents information from an
external world.
This paper proposes an alternative foundation.
A distinction that cannot be restored after perturbation is
operationally indistinguishable from a distinction that never existed.
Observation therefore depends not on reception but on
\emph{restorability}.
We introduce a pre-theoretic notion of \emph{stable observability}
and prove that it coincides exactly with finite-cost restorability
under admissible perturbation---making the title thesis a theorem
rather than a definition.
An observer is then defined as a repair process whose distinction
space is induced by its repair repertoire, and the algebra of repair
spaces is developed: monotonicity, repair equivalence, subadditivity,
the interference functional, the repair pseudometric, and the
Reachability--Restorability inclusion.
The boundary of the observer-relative admissibility manifold $\AO$
is characterised as the divergence locus of minimal repair cost
$\CR$---the Repair Boundary Theorem, grounded in a lemma on lower
semi-continuity of repair cost.
Applications are developed for blindsight and perceptual filling-in,
quantum measurement as repair-or-collapse, and institutional observers.
The framework unifies distinction theory, the admissibility program,
and repair-theoretic accounts of memory and understanding into a
single foundational identification.
\end{abstract}

\vspace{1em}
\hrule
\vspace{1.6em}

% ============================================================
\section{Observation as a Problem}
\label{sec:intro}
% ============================================================

Standard theories of observation begin with reception.
An observer receives signals, records measurements, forms
representations, or updates beliefs~\cite{Marr1982,Shannon1948}.
In every case observation is treated as the acquisition of information
from an external world.
The observer is essentially passive: a container into which the world
deposits its impressions~\cite{Gibson1979}.

This paper begins from a different premise.
A distinction that cannot be restored after perturbation is
operationally indistinguishable from a distinction that never existed.
Observation therefore depends not on reception but on restorability.
What an observer can observe is precisely what it can restore.

We propose that observability and restorability are equivalent notions.
An observer is not fundamentally a receiver of information but a
repair process acting on a substrate of distinctions.
Perception, memory, measurement, and scientific inquiry become special
cases of distinction restoration under perturbation.

The central identification is:
\begin{equation}
  \Obs(O) \;=\; \Rest(O),
  \label{eq:main-id}
\end{equation}
where $\Obs(O)$ is the set of distinctions observable by $O$ and
$\Rest(O)$ is the set of distinctions restorable by $O$.
Everything else in this paper is a consequence of taking
\eqref{eq:main-id} seriously.

The argument proceeds as follows.
Section~\ref{sec:substrate} establishes the mathematical substrate:
state spaces, perturbations, and distinctions.
Section~\ref{sec:repair} defines repair operators, repair paths, and
repair cost.
Section~\ref{sec:restorable} introduces restorable distinctions.
Section~\ref{sec:observer} re-defines the observer as a repair process
whose distinction space is induced by its repair repertoire.
Section~\ref{sec:worlds} studies shared perceptual worlds as
intersections of repair repertoires.
Section~\ref{sec:stable} introduces the pre-theoretic notion of
stable observability and proves that it coincides with restorability,
making \eqref{eq:main-id} a theorem rather than a stipulation.
Section~\ref{sec:structure} develops the internal algebra of repair
spaces: monotonicity, equivalence, subadditivity, the interference
functional, the repair metric, and the reachability--restorability
inclusion.
Section~\ref{sec:boundary} states and proves the Repair Boundary
Theorem, supported by a lemma on lower semi-continuity of repair cost.
Sections~\ref{sec:blindsight}--\ref{sec:institutions} develop
applications.
Section~\ref{sec:conclusion} draws consequences and identifies open
problems.


% ============================================================
\section{Substrates, Perturbations, and Distinctions}
\label{sec:substrate}
% ============================================================

\begin{definition}[Substrate]
  A \emph{substrate} is a triple $S = (\Xsub, \Psub, \Dsub)$ where:
  \begin{itemize}
    \item $\Xsub$ is a topological space of \emph{states};
    \item $\Psub$ is a semigroup of \emph{perturbation operators}
      $p : \Xsub \to \Xsub$ acting continuously on $\Xsub$, closed
      under composition, with identity $\mathrm{id} \in \Psub$;
    \item $\Dsub \subseteq \Xsub \times \Xsub$ is a symmetric,
      irreflexive \emph{distinction relation}: $(x,y) \in \Dsub$
      means states $x$ and $y$ are discriminable.
  \end{itemize}
\end{definition}

Throughout this paper we impose two standing axioms on every substrate.

\begin{axiom}[Locally compact Hausdorff substrate]
  The state space $\Xsub$ is a locally compact Hausdorff topological
  space.
  \label{ax:lch}
\end{axiom}

\begin{axiom}[Closed distinction relation]
  The distinction relation $\Dsub$ is closed in the product topology
  on $\Xsub \times \Xsub$.
  \label{ax:closed-d}
\end{axiom}

Axiom~\ref{ax:lch} provides the compactness and separation needed for
convergence arguments throughout.
Axiom~\ref{ax:closed-d} ensures that limit states of sequences of
discriminable pairs remain discriminable: if $(x_n, y_n) \in \Dsub$
and $(x_n, y_n) \to (x, y)$, then $(x, y) \in \Dsub$.
Without this, the refinement step in Theorem~\ref{thm:stable-obs} can
fail.
Both axioms hold in all concrete examples considered in this paper.

\begin{remark}
  The distinction relation $\Dsub$ need not be transitive and should
  not in general be an equivalence.
  Discrimination is a local, pair-wise affair.
  Two states may be mutually discriminable from a third without being
  mutually discriminable from each other.
\end{remark}

\begin{definition}[Admissible perturbation]
  A perturbation $p \in \Psub$ is \emph{admissible relative to
  distinction $d = (x,y)$} if $p$ maps at least one of $x, y$ to a
  state from which $d$ might fail to hold\,---\,formally, if
  $(p(x), y) \notin \Dsub$ or $(x, p(y)) \notin \Dsub$.
  The set of admissible perturbations for $d$ is denoted $\Psub_d$.
\end{definition}

\begin{definition}[Distinction survival]
  Distinction $d = (x,y)$ \emph{survives perturbation} $p$ if
  $(p(x), p(y)) \in \Dsub$.
  Otherwise $p$ \emph{erases} $d$.
\end{definition}

The central concern of this paper is what happens after erasure.


% ============================================================
\section{Repair Operators and Repair Cost}
\label{sec:repair}
% ============================================================

\begin{definition}[Repair action]
  Given a substrate $S$, a \emph{repair action} is a continuous map
  $r : \Xsub \to \Xsub$ belonging to a designated set
  $\mathcal{R} \subseteq C(\Xsub, \Xsub)$ of available interventions,
  where $C(\Xsub, \Xsub)$ carries the compact-open topology.
  The set $\mathcal{R}$ is the \emph{repair repertoire} of the process
  under consideration.
  The study of controllable interventions on state spaces is treated
  systematically in~\cite{Sontag1998,Kalman1960,Bellman1957}.
\end{definition}

\begin{axiom}[Compact repair repertoire]
  The repair repertoire $\mathcal{R}$ is a sequentially compact subset
  of $C(\Xsub, \Xsub)$ under the compact-open topology.
  Consequently $\mathcal{R}^n$, the product space of repair paths of
  fixed length $n$, is compact by Tychonoff's theorem.
  \label{ax:compact-r}
\end{axiom}

\begin{definition}[Repair path]
  A \emph{repair path} for distinction $d = (x_0, y_0)$ after
  perturbation $p$ is a finite sequence of repair actions
  $\gamma = (r_1, r_2, \ldots, r_n)$ such that, writing
  $x_k = r_k(x_{k-1})$ and $y_k = r_k(y_{k-1})$, we have
  $(x_n, y_n) \in \Dsub$.
  The set of all repair paths of length at most $N$ is denoted
  $\Gamma_d^{p,N}$; the union $\Gamma_d^p = \bigcup_{N \geq 1}
  \Gamma_d^{p,N}$.
\end{definition}

\begin{definition}[Repair cost]
  Let $c : \Xsub \to \mathbb{R}_{\geq 0}$ be a lower semi-continuous
  \emph{local cost density} representing the resource expenditure per
  unit state occupancy.
  Typical interpretations include metabolic energy ($c(x) =$ ATP
  consumption rate), computational effort ($c(x) =$ clock cycles per
  step), or temporal cost ($c(x) = 1$ uniformly).
  The cost of a repair path $\gamma = (r_1, \ldots, r_n)$ is the
  \emph{discrete path sum}
  \[
    \mathrm{cost}(\gamma, x_0)
    \;=\;
    \sum_{k=1}^{n} c\!\left(x_{k-1}\right),
  \]
  where $x_k = r_k(x_{k-1})$.
  When repair actions are parameterised as continuous flows
  $r_k = \Phi_{t_k}$ for $t_k > 0$, this sum extends naturally to
  the line integral $\int_\gamma c\, \mathrm{d}s$; both formulations
  are used below according to context.
\end{definition}

\begin{definition}[Minimal repair cost]
  The \emph{minimal repair cost} of distinction $d$ after perturbation
  $p$, relative to repair repertoire $\mathcal{R}$, is
  \[
    \CR(d, p)
    \;=\;
    \inf_{\gamma \in \Gamma_d^p}
    \mathrm{cost}(\gamma, p(x_0)),
  \]
  setting $\CR(d,p) = +\infty$ when $\Gamma_d^p = \emptyset$.
\end{definition}

\begin{definition}[Worst-case and probabilistic repair cost]
  The \emph{worst-case repair cost} of distinction $d$ is
  \[
    \CR(d)
    \;=\;
    \sup_{p \in \Psub_d}\, \CR(d, p).
  \]
  This is the operationally severe notion: a single catastrophic
  perturbation with $\CR(d,p) = +\infty$ renders $d$ non-restorable
  in the absolute sense.

  For biological and institutional observers, a softer notion is
  often more appropriate.
  Given a probability measure $\mu$ on $\Psub_d$, the
  \emph{probabilistic repair cost} is
  \[
    \CR^\mu(d)
    \;=\;
    \int_{\Psub_d} \CR(d, p)\, \mathrm{d}\mu(p).
  \]
  Absolute restorability $\Rest_\infty(\mathcal{R})$ uses $\CR$;
  probabilistic restorability $\Rest_\mu(\mathcal{R})$ uses
  $\CR^\mu$.
  The former is the object of this paper unless otherwise stated.
\end{definition}


% ============================================================
\section{Restorable Distinctions}
\label{sec:restorable}
% ============================================================

\begin{definition}[Restorability]
  A distinction $d \in \Dsub$ is \emph{restorable by repair repertoire
  $\mathcal{R}$} if $\CR(d) < +\infty$\,---\,that is, if every
  admissible perturbation admits a finite-cost repair path.
  This is \emph{finite-step restorability}; the set is
  \[
    \Restf(\mathcal{R})
    \;=\;
    \bigl\{
      d \in \Dsub \;\big|\; \CR(d) < +\infty
    \bigr\}.
  \]
\end{definition}

The finite-step condition requires a repair path of bounded length.
A weaker but mathematically significant notion permits repair in the
limit of infinitely many steps.

\begin{definition}[Limit restorability]
  A distinction $d \in \Dsub$ is \emph{limit-restorable by
  $\mathcal{R}$} if for every admissible perturbation $p \in \Psub_d$
  and every $\varepsilon > 0$ there exists a (possibly infinite)
  sequence of repair actions $(r_1, r_2, \ldots)$ such that the
  partial orbit $(x_k, y_k)$ satisfies
  \[
    \inf_{(u,v) \in \Dsub}
    \bigl[\|x_k - u\| + \|y_k - v\|\bigr]
    \;\to\; 0
    \quad\text{as } k \to \infty,
  \]
  with $\sum_{k=1}^\infty c(x_{k-1}) < +\infty$.
  The set of limit-restorable distinctions is $\Reston(\mathcal{R})$.
\end{definition}

\begin{proposition}[Inclusion of restorability classes]
  $\Restf(\mathcal{R}) \subseteq \Reston(\mathcal{R})$.
  The inclusion is in general strict.
  \label{prop:class-inclusion}
\end{proposition}

\begin{proof}
  Any finite repair path is a finite sequence; the partial orbit
  reaches $\Dsub$ in finitely many steps, so the infimum condition
  holds in the limit.
  Strictness: consider a substrate in which every perturbation of $d$
  can be countered by successively halving the remaining distance to
  $\Dsub$, with costs $1/2^k$; the total cost $\sum 1/2^k = 1$ is
  finite, but no single finite path lands in the closed set $\Dsub$
  unless the limit state is taken.
  Under Axiom~\ref{ax:closed-d} the limit is in $\Dsub$, confirming
  limit-restorability; but if no path of fixed finite length suffices,
  the distinction is in $\Reston \setminus \Restf$.
\end{proof}

\begin{remark}
  The two classes separate different physical regimes.
  Biological repair (neural recovery, wound healing) and mathematical
  proof search aspire to $\Restf$: a bounded sequence of discrete
  interventions must suffice.
  Iterative scientific refinement, asymptotic stabilisation of
  institutions, and language reconstruction may be modelled more
  faithfully by $\Reston$: convergence is guaranteed but no finite
  stopping point is specified.
  Courts occupy an intermediate position, tolerating appellate
  iteration but demanding termination.
  The Stable Observability Theorem (Theorem~\ref{thm:stable-obs}) is
  a theorem about $\Restf$; a parallel theorem for $\Reston$ under
  weaker topological assumptions is an open problem
  (Problem~1 in Section~\ref{sec:conclusion}).
\end{remark}

Throughout Sections~\ref{sec:observer}--\ref{sec:boundary},
$\Rest(\mathcal{R})$ denotes $\Restf(\mathcal{R})$ unless otherwise
stated.

\begin{proposition}[Restorability is closed under finite composition]
  If $d_1, d_2 \in \Rest(\mathcal{R})$ and the restoration of $d_2$
  does not interfere with the restored state of $d_1$ (formally: the
  repair paths commute on the relevant component of $\Xsub$), then
  the compound distinction $d_1 \wedge d_2$ is also in
  $\Rest(\mathcal{R})$.
  \label{prop:closure}
\end{proposition}

\begin{proof}
  % SKELETON: concatenate repair paths; bound cost by sum.
  Let $\gamma_1 \in \Gamma_{d_1}^p$ and $\gamma_2 \in \Gamma_{d_2}^p$
  be finite-cost repair paths.
  Under the commutativity assumption, $\gamma_1 \cdot \gamma_2$
  (concatenation) restores both distinctions with cost at most
  $\mathrm{cost}(\gamma_1) + \mathrm{cost}(\gamma_2) < +\infty$.
\end{proof}

\begin{remark}
  The commutativity assumption in Proposition~\ref{prop:closure} is
  non-trivial and fails in general.
  Section~\ref{sec:blindsight} exhibits cases where restoring one
  distinction actively interferes with another,
  producing \emph{repair conflict}.
\end{remark}


% ============================================================
\section{Observers as Repair Processes}
\label{sec:observer}
% ============================================================

The standard move is to define an observer as a pair $(O, \DO)$ of a
process and a set of distinctions it tracks~\cite{Ashby1956,ConantAshby1970}.
We invert this dependency.

\begin{definition}[Observer]
  An \emph{observer} is a repair process $O = (S, \mathcal{R}_O)$
  consisting of a substrate $S$ and a repair repertoire
  $\mathcal{R}_O$.
  The \emph{observation space} of $O$ is the induced set
  \[
    \DO \;=\; \Rest(\mathcal{R}_O).
  \]
\end{definition}

The distinction space is not given in advance and then repaired.
It is the repair repertoire that determines which distinctions exist
for this observer.
An observer with a larger repertoire inhabits a richer distinction
space, not because it receives more information but because it can
restore more.

\begin{definition}[Observability]
  A distinction $d$ is \emph{observable by $O$} if and only if
  $d \in \Rest(\mathcal{R}_O)$.
  In symbols:
  \[
    \Obs(O) \;=\; \Rest(\mathcal{R}_O) \;=\; \DO.
  \]
\end{definition}

This is the formal content of \eqref{eq:main-id}.
Observability just is restorability.


% ============================================================
\section{Perceptual Worlds and Observer Agreement}
\label{sec:worlds}
% ============================================================

\begin{definition}[Shared perceptual world]
  Two observers $O_1$ and $O_2$ \emph{share a perceptual world} to
  the extent that their observation spaces overlap:
  \[
    W(O_1, O_2)
    \;=\;
    \Obs(O_1) \cap \Obs(O_2)
    \;=\;
    \Rest(\mathcal{R}_{O_1}) \cap \Rest(\mathcal{R}_{O_2}).
  \]
\end{definition}

\begin{theorem}[Shared perception is shared repair]
  Two observers share a perceptual world if and only if they share
  repair capacity over that world:
  \[
    W(O_1, O_2)
    \;=\;
    \Rest(\mathcal{R}_{O_1}) \cap \Rest(\mathcal{R}_{O_2}).
  \]
  In particular, two observers with disjoint repair repertoires
  inhabit disjoint perceptual worlds even when embedded in the same
  physical substrate.
  \label{thm:shared-world}
\end{theorem}

\begin{proof}
  Immediate from Definition~\ref{sec:observer} and the definition of
  $W(O_1, O_2)$.
\end{proof}

\begin{corollary}[Observer agreement]
  Two observers \emph{agree} on distinction $d$ if and only if
  $d \in W(O_1, O_2)$, i.e., both can restore $d$ after any
  admissible perturbation.
  Agreement on signals is neither necessary nor sufficient;
  what matters is agreement on repairability.
\end{corollary}

\begin{example}[Species-relative perception]
  A bat and a human occupy the same physical environment.
  Their repair repertoires differ substantially in the echoic
  dimension.
  Theorem~\ref{thm:shared-world} predicts that the
  spatial distinctions available to the bat via echolocation are
  not observable by the human, not because the signals are absent but
  because the human lacks repair operators acting on that channel.
\end{example}


% ============================================================
\section{Stable Observability}
\label{sec:stable}
% ============================================================

The identification $\Obs(O) = \Rest(O)$ as stated in
Section~\ref{sec:observer} is definitional: observability was
\emph{defined} as restorability.
A genuine theorem requires a pre-theoretic notion of observability
that is independent of the repair apparatus, so that the equivalence
can be established rather than assumed.

We supply one now.

\begin{definition}[Stable observability]
  Fix an observer $O = (S, \mathcal{R}_O)$ and a threshold
  $\varepsilon > 0$.
  A distinction $d = (x, y)$ is \emph{stably observable by $O$} if
  for every $\delta > 0$ and every admissible perturbation
  $p \in \Psub_d$, there exists a finite sequence of repair actions
  $(r_1, \ldots, r_n) \in \mathcal{R}_O^*$ such that the resulting
  states $(x', y')$ satisfy
  \[
    \inf_{(u,v) \in \Dsub} \bigl[\,
      \|x' - u\| + \|y' - v\|
    \,\bigr] < \delta,
  \]
  and the total cost of the sequence is bounded independently of
  $\delta$.
  The set of stably observable distinctions is denoted
  $\ObsStab(O)$.
\end{definition}

Informally: $d$ is stably observable if, no matter how it is
perturbed, $O$ can return to a state of discrimination within finite
cost.
The definition is pre-theoretic in the sense that it makes no
reference to repair paths as formal objects; it speaks only of the
capacity to discriminate after perturbation.

\begin{theorem}[Stable Observability Theorem]
  Under Axioms~\ref{ax:lch}, \ref{ax:closed-d}, and~\ref{ax:compact-r},
  and assuming $c : \Xsub \to \mathbb{R}_{\geq 0}$ is lower
  semi-continuous,
  \[
    \ObsStab(O) \;=\; \Rest(\mathcal{R}_O).
  \]
  \label{thm:stable-obs}
\end{theorem}

\begin{proof}
  $(\supseteq)$\;
  Suppose $d \in \Rest(\mathcal{R}_O)$.
  Then $\CR(d) < +\infty$, so every admissible perturbation
  admits a finite-cost repair path $\gamma^* \in \Gamma_d^p$ that
  lands in $\Dsub$.
  The discrimination condition in the definition of stable
  observability is satisfied with $\delta = 0$ (strictly inside
  $\Dsub$), and the cost bound $C = \CR(d)$ is independent of
  $\delta$.
  Thus $d \in \ObsStab(O)$.

  $(\subseteq)$\;
  Suppose $d = (x_0, y_0) \in \ObsStab(O)$.
  Fix an admissible perturbation $p \in \Psub_d$.
  By stable observability, for each $m \geq 1$ there exists a finite
  repair sequence $\gamma_m = (r_{1,m}, \ldots, r_{n_m, m})$ of
  total cost $\mathrm{cost}(\gamma_m) \leq C_p < +\infty$ (with $C_p$
  independent of $m$) such that the terminal states
  $(x'_m, y'_m)$ satisfy
  \[
    \inf_{(u,v) \in \Dsub}
    \bigl[\|x'_m - u\| + \|y'_m - v\|\bigr]
    < \tfrac{1}{m}.
  \]

  \textbf{Bounding the path length.}\;
  Since $c$ is lower semi-continuous and positive, and $\Xsub$ is
  locally compact, there exists a lower bound $c_0 > 0$ on any compact
  neighbourhood of the orbit.
  Each repair step contributes at least $c_0$ to the cost, so
  $n_m \leq C_p / c_0 =: N$ uniformly in $m$.
  We may therefore assume all paths have the same length $n \leq N$
  (padding with the identity if needed).

  \textbf{Extracting a convergent subsequence.}\;
  By Axiom~\ref{ax:compact-r}, $\mathcal{R}_O$ is sequentially compact
  under the compact-open topology, so $\mathcal{R}_O^n$ is compact by
  Tychonoff.
  The sequence of tuples $(\gamma_m)_{m \geq 1}$ in $\mathcal{R}_O^n$
  admits a convergent subsequence
  $(\gamma_{m_j}) \to \gamma^* = (r_1^*, \ldots, r_n^*)$.
  By joint continuity of evaluation in the compact-open topology on
  locally compact Hausdorff spaces (Axiom~\ref{ax:lch}), the terminal
  states converge:
  $(x'_{m_j}, y'_{m_j}) \to (x'^*, y'^*)$.

  \textbf{Landing in $\Dsub$.}\;
  For each $j$, the terminal states satisfy
  $\inf_{(u,v) \in \Dsub}[\|x'_{m_j} - u\| + \|y'_{m_j} - v\|]
  < 1/{m_j}$,
  so $(x'_{m_j}, y'_{m_j})$ is within $1/{m_j}$ of $\Dsub$.
  In the limit, $(x'^*, y'^*) \in \overline{\Dsub} = \Dsub$ by
  Axiom~\ref{ax:closed-d}.
  Therefore $\gamma^* \in \Gamma_d^p$ is a valid repair path.

  \textbf{Cost bound.}\;
  The total cost functional
  $\mathrm{cost}(\gamma, x) = \sum_{k=1}^n c(x_{k-1})$ is a finite
  sum of lower semi-continuous functions under the compact-open
  topology and is therefore lower semi-continuous itself.
  Thus
  \[
    \mathrm{cost}(\gamma^*) \;\leq\; \liminf_{j} \mathrm{cost}(\gamma_{m_j})
    \;\leq\; C_p \;<\; +\infty.
  \]
  Since this holds for every admissible perturbation $p$,
  $\CR(d) \leq C_p < +\infty$, so $d \in \Rest(\mathcal{R}_O)$.
\end{proof}

\begin{remark}
  Theorem~\ref{thm:stable-obs} is the paper's load-bearing result.
  It establishes that the title thesis---observability is
  restorability---is not a definitional convenience but a substantive
  claim: the pre-theoretic capacity to maintain discrimination under
  perturbation is exactly captured by the formal notion of finite-cost
  repair.
  Every subsequent theorem is a consequence of taking this equivalence
  seriously.
\end{remark}


% ============================================================
\section{The Algebra of Repair Spaces}
\label{sec:structure}
% ============================================================

We now develop the internal mathematics of repair spaces before
connecting them to the admissibility manifold.
The results of this section do not depend on
Theorem~\ref{thm:stable-obs} and hold for any repair repertoire.

\subsection{Monotonicity}

\begin{theorem}[Repair Monotonicity]
  Let $\mathcal{R}_1 \subseteq \mathcal{R}_2$ be repair repertoires
  acting on the same substrate $S$.
  Then
  \[
    \Rest(\mathcal{R}_1)
    \;\subseteq\;
    \Rest(\mathcal{R}_2).
  \]
  \label{thm:monotonicity}
\end{theorem}

\begin{proof}
  If $d \in \Rest(\mathcal{R}_1)$, every admissible perturbation
  of $d$ admits a finite-cost repair path composed entirely of
  elements of $\mathcal{R}_1$.
  Since $\mathcal{R}_1 \subseteq \mathcal{R}_2$, every such path is
  available in $\mathcal{R}_2$, so $\CR[\mathcal{R}_2](d) \leq
  \CR[\mathcal{R}_1](d) < +\infty$.
  Thus $d \in \Rest(\mathcal{R}_2)$.
\end{proof}

\begin{corollary}[Learning enlarges worlds]
  Whenever $\mathcal{R}_{O_1} \subseteq \mathcal{R}_{O_2}$,
  \[
    \Obs(O_1) \;\subseteq\; \Obs(O_2).
  \]
  Acquiring new repair operators strictly enlarges the set of
  observable distinctions.
  This gives a precise sense in which learning enlarges the world an
  observer inhabits: not by receiving more signals, but by gaining
  more repair capacity.
\end{corollary}

\subsection{Observer Equivalence}

\begin{definition}[Repair equivalence]
  Observers $O_1$ and $O_2$ are \emph{repair-equivalent},
  written $O_1 \sim_R O_2$, if
  $\Rest(\mathcal{R}_{O_1}) = \Rest(\mathcal{R}_{O_2})$.
\end{definition}

\begin{theorem}[Repair Equivalence Theorem]
  Repair-equivalent observers are observationally indistinguishable:
  $O_1 \sim_R O_2$ implies $\Obs(O_1) = \Obs(O_2)$.
  \label{thm:equivalence}
\end{theorem}

\begin{proof}
  Immediate from the identification $\Obs(O) = \Rest(\mathcal{R}_O)$.
\end{proof}

\begin{remark}
  Theorem~\ref{thm:equivalence} is a repair-theoretic analogue of
  bisimulation in process calculi.
  Two observers may have entirely different internal architectures
  and yet be observationally equivalent if their repair repertoires
  generate the same restorable set.
  This suggests that the natural equivalence class for observer
  identity is repair capacity, not physical constitution---a
  conclusion consonant with the autopoietic account of
  identity~\cite{MaturanaVarela1980,VarelaThompsonRosch1991} but
  grounded here in the repair structure rather than in biological
  self-production.
\end{remark}

\subsection{Subadditivity and the Interference Functional}

\begin{theorem}[Subadditivity of Repair Cost]
  Suppose distinctions $d_1$ and $d_2$ admit commuting repair paths
  (as in Proposition~\ref{prop:closure}).
  Then
  \[
    \CR(d_1 \wedge d_2)
    \;\leq\;
    \CR(d_1) + \CR(d_2).
  \]
  \label{thm:subadditive}
\end{theorem}

\begin{proof}
  Let $\gamma_1^* \in \Gamma_{d_1}^p$ and $\gamma_2^* \in
  \Gamma_{d_2}^p$ be minimal repair paths for each distinction
  separately.
  Under commutativity, the concatenation
  $\gamma_1^* \cdot \gamma_2^*$ is a valid repair path for
  $d_1 \wedge d_2$ with cost at most
  $\mathrm{cost}(\gamma_1^*) + \mathrm{cost}(\gamma_2^*)$.
  Taking infima on each side yields the stated inequality.
\end{proof}

\begin{definition}[Repair interference functional]
  For distinctions $d_1, d_2 \in \Rest(\mathcal{R})$, define
  \[
    I(d_1, d_2)
    \;=\;
    \CR(d_1 \wedge d_2) - \CR(d_1) - \CR(d_2).
  \]
  \label{def:interference}
\end{definition}

By Theorem~\ref{thm:subadditive}, $I(d_1, d_2) \leq 0$ when repair
paths commute.
Three regimes are distinguished:

\begin{itemize}
  \item $I(d_1, d_2) > 0$: \emph{repair conflict}.
    Restoring both distinctions jointly costs more than restoring
    each separately.
    This is the formal signature of perceptual trade-offs:
    attending to one distinction impedes maintenance of another.

  \item $I(d_1, d_2) = 0$: \emph{repair independence}.
    The distinctions are orthogonal in repair space.

  \item $I(d_1, d_2) < 0$: \emph{repair synergy}.
    Joint restoration is cheaper than independent restoration.
    This captures perceptual binding: some distinctions are
    more efficiently maintained together than apart, which
    may underlie Gestalt grouping effects and multi-modal
    perceptual integration.
\end{itemize}

When the restorable distinction space is finite or countable,
the pairwise values $I(d_i, d_j)$ assemble into a matrix.

\begin{definition}[Interference matrix]
  Let $\{d_1, \ldots, d_n\} \subseteq \Rest(\mathcal{R})$ be a finite
  collection of distinctions.
  The \emph{interference matrix} is
  \[
    \IntMat \;=\; \bigl(I(d_i, d_j)\bigr)_{i,j=1}^n \;\in\; \mathbb{R}^{n \times n}.
  \]
  \label{def:int-matrix}
\end{definition}

\begin{proposition}[Spectral interpretation of the interference matrix]
  $\IntMat$ is symmetric.
  Its spectrum $\spec(\IntMat)$ carries the following interpretation:
  \begin{enumerate}
    \item \emph{Large positive eigenvalues}: the corresponding
      eigenvectors identify \emph{repair bottleneck clusters}---groups
      of distinctions that mutually impede each other's restoration.
      These correspond to attentional capacity limits in biological
      observers and bandwidth constraints in institutional ones.
    \item \emph{Large negative eigenvalues}: the eigenvectors identify
      \emph{synergistic bundles}---groups of distinctions that are
      substantially cheaper to restore jointly.
      These correspond to Gestalt perceptual groupings, co-regulated
      gene networks, and inter-dependent institutional procedures.
    \item \emph{Near-zero eigenvalues}: the associated distinctions are
      approximately independent in repair space---neither conflicting
      nor synergistic.
  \end{enumerate}
  The Frobenius norm $\|\IntMat\|_F$ quantifies the total interaction
  strength; $\|\IntMat\|_F = 0$ iff all distinctions in the collection
  are repair-independent.
  \label{prop:spectrum}
\end{proposition}

\begin{proof}
  Symmetry: $I(d_i, d_j) = I(d_j, d_i)$ since $d_i \wedge d_j =
  d_j \wedge d_i$ and $\CR$ is a function of the distinction alone.
  The spectral interpretations follow from standard spectral geometry:
  the quadratic form $v^\top \IntMat v$ for a unit vector $v$
  represents the weighted interference over the combination
  $\sum_i v_i d_i$, with positive values indicating net conflict and
  negative values indicating net synergy.
\end{proof}

\begin{remark}
  The interference matrix has the algebraic structure of a signed
  Gram matrix over the distinction space, and its spectral theory
  connects the repair framework to random matrix theory and
  the theory of signed networks.
  For an observer with a large restorable distinction space, the
  density of states of $\IntMat$ describes the global interaction
  structure of perception.
  A repair bottleneck corresponds to a rank-one positive perturbation
  of $\IntMat$; the eigenvalue gap around zero measures how close
  the observer is to a globally synergistic or globally conflicted
  regime.
  Ray~Barman et al.~\cite{RayBarman2026} demonstrate empirically that
  interference among competing memories---rather than temporal decay
  alone---drives power-law forgetting in high-dimensional retrieval
  systems, and that vulnerability to interference is governed by
  \emph{effective} rather than nominal dimensionality.
  The interference matrix provides a formal object in which these
  empirical phenomena have a natural home: the effective dimensionality
  of a retrieval space corresponds to the rank of $\IntMat$, and the
  concentration of variance into few dimensions corresponds to a
  low-rank positive-semidefinite structure with a small number of
  large positive eigenvalues.
\end{remark}

\subsection{The Repair Pseudometric}

\begin{definition}[Repair pseudometric]
  For distinctions $d_i, d_j \in \Rest(\mathcal{R})$, define
  \[
    d_R(d_i, d_j)
    \;=\;
    \inf_{\gamma : d_i \leadsto d_j}
    \int_{\gamma} c(x)\, \mathrm{d}s,
  \]
  where the infimum is taken over all repair trajectories in
  $\Xsub$ that transform the substrate state associated with
  $d_i$ into one from which $d_j$ is restorable.
  The construction is formally analogous to the Kantorovich--Wasserstein
  distance on probability measures~\cite{Villani2009}, with repair
  trajectories playing the role of transport plans.
\end{definition}

\begin{proposition}[Repair pseudometric properties]
  The function $d_R$ satisfies:
  \begin{enumerate}
    \item $d_R(d_i, d_i) = 0$ for all $d_i$;
    \item $d_R(d_i, d_k) \leq d_R(d_i, d_j) + d_R(d_j, d_k)$
      (triangle inequality, by path concatenation);
    \item If repair trajectories are reversible and $c$ is symmetric
      under reversal, then $d_R(d_i, d_j) = d_R(d_j, d_i)$ and
      $d_R$ is a genuine metric on $\Rest(\mathcal{R})$.
  \end{enumerate}
  \label{prop:metric}
\end{proposition}

\begin{remark}
  Irreversible perturbations---cell death, archive destruction,
  data erasure---produce cases where $d_R(d_i, d_j) \neq d_R(d_j, d_i)$.
  The asymmetric pseudometric is therefore the primary object; the
  metric is a special case arising under reversibility.
  The directed structure of irreversible repair spaces may connect to
  the time-asymmetric picture developed in the RSVP framework.
\end{remark}

\subsection{Reachability and Restorability}

The following theorem is the formal weld between the repair-observer
framework and the admissibility program.

\begin{theorem}[Reachability--Restorability Inclusion]
  Let $\AO$ be the observer-relative admissibility manifold
  (Definition of Section~\ref{sec:boundary}) and let
  $\Reach(\AO)$ denote the set of distinctions reachable
  within $\AO$ under the dynamics of $S$.
  Then
  \[
    \Rest(\mathcal{R}_O)
    \;\subseteq\;
    \Reach(\AO).
  \]
  Every observable distinction lies within the observer's reachable
  admissibility region.
  \label{thm:reach-rest}
\end{theorem}

\begin{proof}
  Suppose $d \in \Rest(\mathcal{R}_O)$.
  Then $\CR(d) < +\infty$, so there exists a finite-cost repair path
  $\gamma$ that restores $d$ from any admissible perturbation.
  This path lies within $\Xsub$ and terminates in a state where $d$
  is discriminable.
  By definition, the set of states from which $d$ remains restorable
  is contained in $\AO$ (since those states are precisely those from
  which the observer can maintain its distinctions).
  The path $\gamma$ therefore connects the perturbed state to a point
  in $\AO$, establishing that $d$ is reachable within $\AO$.
\end{proof}

\begin{remark}
  The inclusion $\Rest(\mathcal{R}_O) \subseteq \Reach(\AO)$ is in
  general strict.
  Equality would require that every distinction reachable within the
  admissibility manifold is also restorable by the repair repertoire,
  which fails when the manifold contains states that are geometrically
  accessible but dynamically blocked.
  Whether equality holds under compactness and convexity assumptions
  is an open question (Problem~5 in Section~\ref{sec:conclusion}).
\end{remark}

\subsection{Control--Geometry Duality}

The Reachability--Restorability Inclusion reveals a structural duality
running through the entire framework.
On one side stands the \emph{control} perspective: the repair
repertoire $\mathcal{R}_O$ generates a dynamics on $\Xsub$, and
questions of repair are questions of controllability.
On the other side stands the \emph{geometry} perspective: the restorable
distinction space $\Rest(\mathcal{R}_O)$ inherits a metric structure
from repair cost, and questions of observation are questions of
geometry.

\begin{definition}[Control structure and induced geometry]
  The \emph{control structure} of an observer is the triple
  $\mathfrak{C}(\mathcal{R}_O) = (\Xsub, \mathcal{R}_O, c)$:
  a state space, a set of admissible interventions, and a cost
  density.
  The \emph{induced geometry} is the pair
  $\mathfrak{G}(\Rest(\mathcal{R}_O)) = (\Rest(\mathcal{R}_O), d_R)$:
  the restorable distinction space equipped with the repair
  pseudometric.
\end{definition}

The duality takes the form of a correspondence:

\begin{center}
\renewcommand{\arraystretch}{1.4}
\begin{tabular}{l@{\quad$\longleftrightarrow$\quad}l}
  Reachability in $\mathfrak{C}$ & Existence of a repair trajectory \\
  Restorability in $\mathfrak{G}$ & Finite cost of a repair trajectory \\
  Admissibility boundary $\bA_O$ & Divergence locus of $\CR$ \\
  Local controllability & Open geodesic balls in $d_R$ \\
  Topological controllability & Dense repair orbits in $\Rest(\mathcal{R}_O)$ \\
\end{tabular}
\end{center}

\begin{theorem}[Control--Geometry Collapse]
  Let $\Xsub \subseteq \mathbb{R}^m$ be compact and convex and let
  the repair repertoire $\mathcal{R}_O$ consist of continuous maps
  forming a locally state-controllable system on the interior of
  $\AO$.
  Then
  \[
    \Rest(\mathcal{R}_O) \;=\; \Reach(\AO).
  \]
  \label{thm:cg-collapse}
\end{theorem}

\begin{proof}
  $(\subseteq)$\; Theorem~\ref{thm:reach-rest}.

  $(\supseteq)$\;
  Suppose $d \in \Reach(\AO)$.
  By local state-controllability, the reachable set from any interior
  point of $\AO$ is open.
  Any continuous path in $\AO$ from the perturbed state to a state
  where $d$ is discriminable can therefore be covered by finitely many
  open controllable neighbourhoods (by compactness of $\AO$).
  Each neighbourhood transition is realised by a finite sequence of
  repair operators in $\mathcal{R}_O$.
  The composite path has total cost bounded by $N \cdot M$ where
  $N$ is the number of transitions and $M = \max_{x \in \Xsub} c(x)
  < +\infty$ (since $c$ is lower semi-continuous on a compact space).
  Thus $\CR(d) \leq N \cdot M < +\infty$, so $d \in \Rest(\mathcal{R}_O)$.
\end{proof}

\begin{remark}
  Theorem~\ref{thm:cg-collapse} is a resolution of Problem~5 under
  explicit hypotheses.
  The control-theoretic condition (local state-controllability) is the
  precise content of the intuition that ``every geometrically accessible
  state is dynamically reachable.''
  When $\mathcal{R}_O$ generates a dense subgroup of the homeomorphism
  group of $\Xsub$, controllability holds automatically and the
  equality follows.
  The Repair Boundary Theorem is then a corollary of the duality:
  the boundary $\bA_O$ is simultaneously the set of states from which
  the geometry runs out (repair cost diverges) and the set of states
  from which the control system can no longer steer.
\end{remark}




We characterise the boundary of the observer-relative admissibility
manifold in terms of diverging repair cost.
Theorem~\ref{thm:reach-rest} established that every restorable
distinction lies within the reachable admissibility region;
the Repair Boundary Theorem pinpoints where that region ends.
It requires a preparatory lemma.

\begin{lemma}[Lower semi-continuity of repair cost]
  Let $\Xsub$ be locally compact (Axiom~\ref{ax:lch}), let the local
  cost density $c : \Xsub \to \mathbb{R}_{\geq 0}$ be lower
  semi-continuous~\cite{RockafellarWets1998}, and let the admissible
  repair paths form a family closed under uniform limits on compact
  intervals.
  Then the minimal repair cost
  \[
    \CR(d, x)
    \;=\;
    \inf_{\gamma \in \Gamma_d(x)}
    \int_{\gamma} c\, \mathrm{d}s
  \]
  is lower semi-continuous in $x$.
  \label{lem:lsc}
\end{lemma}

\begin{proof}
  Let $(x_n)$ be a sequence in $\Xsub$ converging to $x$.
  We show $\liminf_{n \to \infty} \CR(d, x_n) \geq \CR(d, x)$.

  For each $n$, let $\gamma_n \in \Gamma_d(x_n)$ be a near-optimal
  repair path with
  $\mathrm{cost}(\gamma_n) \leq \CR(d, x_n) + 1/n$.
  Restricting to a compact neighbourhood $K$ of $x$ and using
  local compactness, the paths $(\gamma_n)$ are equibounded in cost.
  By the assumed closure under uniform limits on compact intervals,
  a subsequence $(\gamma_{n_k})$ converges uniformly to a limit path
  $\gamma^* \in \Gamma_d(x)$.

  By lower semi-continuity of $c$ and Fatou's lemma applied to the
  path integrals,
  \[
    \int_{\gamma^*} c\, \mathrm{d}s
    \;\leq\;
    \liminf_{k \to \infty}
    \int_{\gamma_{n_k}} c\, \mathrm{d}s
    \;\leq\;
    \liminf_{n \to \infty} \CR(d, x_n).
  \]
  Since $\gamma^* \in \Gamma_d(x)$, we have
  $\CR(d, x) \leq \int_{\gamma^*} c\, \mathrm{d}s$,
  completing the proof.
\end{proof}

\begin{definition}[Observer-relative admissibility manifold]
  Let $O = (S, \mathcal{R}_O)$ be an observer.
  The \emph{observer-relative admissibility manifold} $\AO$ is the
  closure of the set of states from which all distinctions in $\DO$
  remain restorable:
  \[
    \AO
    \;=\;
    \overline{\bigl\{
      x \in \Xsub
      \;\big|\;
      \forall\, d \in \DO,\;
      \CR(d) < +\infty \text{ from } x
    \bigr\}}.
  \]
\end{definition}

\begin{theorem}[Repair Boundary Theorem]
  A state $x$ lies on the boundary of the observer-relative
  admissibility manifold if and only if the minimal repair cost of at
  least one distinction in $\DO$ diverges as the system approaches $x$:
  \[
    x \in \bA_O
    \quad\Longleftrightarrow\quad
    \exists\, d \in \DO : \CR(d) \to +\infty
    \text{ as the system state approaches } x.
  \]
  \label{thm:repair-boundary}
\end{theorem}

\begin{proof}
  $(\Rightarrow)$\; Suppose $x \in \bA_O$.
  Then $x$ is a limit point of $\AO$ and also of its complement.
  In any neighbourhood $U \ni x$ there exist states from which some
  distinction $d \in \DO$ is not finitely restorable.
  By Lemma~\ref{lem:lsc}, $\CR(d, \cdot)$ is lower semi-continuous;
  hence $\CR(d)$ must become unbounded along any sequence approaching
  $x$ from outside $\AO$, establishing $\CR(d) \to +\infty$.

  $(\Leftarrow)$\; If $\CR(d) \to +\infty$ as the state approaches
  $x$, then in any neighbourhood of $x$ there are states from which
  $d$ is not finitely restorable.
  Hence $x \notin \mathrm{int}(\AO)$, so $x \in \bA_O$.
\end{proof}

\begin{remark}
  Theorem~\ref{thm:repair-boundary} identifies the admissibility
  boundary not as an externally imposed constraint but as an emergent
  consequence of the repair structure.
  The boundary is where the observer's world ends, in the precise sense
  that no further repair is possible.
\end{remark}

\begin{corollary}[Admissibility as repair capacity]
  The interior of $\AO$ is exactly the set of states from which the
  observer can maintain all its distinctions under all admissible
  perturbations.
  Moving toward $\bA_O$ is operationally equivalent to losing
  observational capacity.
\end{corollary}


% ============================================================
\section{Blindsight and Perceptual Filling-In}
\label{sec:blindsight}
% ============================================================

The repair-observer framework generates a clean account of
\emph{blindsight}: the phenomenon in which subjects with cortical
damage respond discriminatively to visual stimuli they cannot
consciously report~\cite{Weiskrantz1986}.
The division of visual processing into action-guiding and
perception-reporting streams~\cite{MilnerGoodale2006} maps directly
onto the sub-report/report-level restorability split defined below.

In the standard receptive model, blindsight is paradoxical: how can
information be used if it is not received at the reporting level?

Under the repair model the paradox dissolves.
Let the observer $O$ decompose into repair sub-processes operating at
different substrate levels: $O = O_{\mathrm{sub}} \cup O_{\mathrm{rep}}$,
where $O_{\mathrm{sub}}$ acts below the reporting interface and
$O_{\mathrm{rep}}$ acts at the interface itself.

\begin{definition}[Sub-report restorability]
  A distinction $d$ is \emph{sub-report restorable} if
  $d \in \Rest(\mathcal{R}_{O_{\mathrm{sub}}})$ but
  $d \notin \Rest(\mathcal{R}_{O_{\mathrm{rep}}})$.
\end{definition}

Blindsight is then the condition in which the distinction is
sub-report restorable: the sub-cortical repair process maintains the
distinction, but the cortical repair process (whose outputs constitute
verbal report) does not have access to a repair path.
The subject \emph{observes} in the sub-report sense while failing to
\emph{observe} in the report sense.

\begin{remark}
  This account predicts that blindsight subjects should show
  above-chance performance on forced-choice tasks (sub-report
  restoration succeeds) while failing to initiate voluntary report
  (report-level restoration is unavailable).
  This is the observed pattern.
\end{remark}

Perceptual filling-in admits a parallel treatment.
The visual system restores distinctions at the boundary of the blind
spot not by receiving signals from that region but by extending repair
trajectories from surrounding states.
The filled-in percept is a repair product, not a received datum.


% ============================================================
\section{Quantum Measurement as Repair or Collapse}
\label{sec:quantum}
% ============================================================

Standard quantum measurement theory distinguishes pre-measurement
superposition from post-measurement eigenstate.
The \emph{measurement problem} asks what physical process constitutes
the collapse.

The repair-observer framework does not solve the measurement problem
but reframes it productively.
The present account is complementary to relational quantum
mechanics~\cite{Rovelli1996}, which treats measurement outcomes as
relative to observers, and to the decoherence program~\cite{Zurek2003},
which identifies the mechanism by which classical pointer states
become stable.
Under the repair model, decoherence is the process by which the
apparatus acquires repair operators capable of acting on classical
pointer distinctions.

Let the substrate $S$ be a quantum system with Hilbert space
$\mathcal{H}$ and let the distinction relation $\Dsub$ be defined on
density operators: $(\rho_1, \rho_2) \in \Dsub$ iff
$\|\rho_1 - \rho_2\|_1 > \varepsilon$ for some threshold
$\varepsilon > 0$.

A perturbation $p$ that places the system in a superposition
$|\psi\rangle = \alpha|0\rangle + \beta|1\rangle$
erases the distinction between $|0\rangle$ and $|1\rangle$ at the
classical level (neither eigenstate is discriminated from the other
in $\mathcal{H}$ under $p$).

Measurement is then a repair event: either the apparatus-observer
$O_{\mathrm{app}}$ possesses a repair operator that drives the state
to a definite eigenstate (collapse as successful restoration), or it
does not (the distinction remains unrestorable at the classical
interface).

\begin{proposition}[Measurement as repair completion]
  A measurement interaction constitutes observation of distinction
  $d = (|0\rangle, |1\rangle)$ by apparatus $O_{\mathrm{app}}$ if and
  only if $O_{\mathrm{app}}$ possesses a repair path that takes the
  post-perturbation superposition to a state in which $d$ is
  discriminable at the pointer basis.
  \label{prop:measurement}
\end{proposition}

\begin{remark}
  This connects to the quadratic gravity interface picture developed
  elsewhere in the Flyxion program: the interface at which quantum
  distinctions become classically available is the interface at which
  repair operators acting on the classical substrate can restore those
  distinctions.
  Decoherence is the mechanism by which this interface becomes
  accessible.
  %% EXPAND: connection to RSVP interface papers.
\end{remark}


% ============================================================
\section{Institutional Observers}
\label{sec:institutions}
% ============================================================

Individual biological observers are a special case.
The framework is equally applicable to any process with a repair
repertoire.

\begin{definition}[Institutional observer]
  An \emph{institutional observer} is an observer
  $O_{\mathrm{inst}} = (S, \mathcal{R}_{\mathrm{inst}})$ in which the
  substrate $S$ is a social or informational system and the repair
  repertoire $\mathcal{R}_{\mathrm{inst}}$ consists of codified
  institutional procedures.
\end{definition}

\begin{example}[Legal systems]
  A court is an institutional observer whose substrate is the space of
  legally characterised states and whose repair repertoire consists of
  evidentiary and procedural rules.
  A distinction $d$ between \emph{legal} and \emph{illegal} conduct is
  observable by the court if and only if the evidentiary procedures
  can restore that distinction after the perturbation of competing
  characterisations.

  When courts collapse or procedures are suspended, the distinction
  does not become \emph{false}: it becomes \emph{unrestorable}.
  This is the precise sense in which institutional decay is
  observational decay.
\end{example}

\begin{example}[Archives and libraries]
  An archive restores distinctions across time by maintaining repair
  paths from present queries to past states.
  A destroyed archive does not falsify historical distinctions; it
  removes the repair operators by which those distinctions could be
  recovered.
\end{example}

\begin{example}[Scientific laboratories]
  A laboratory is an institutional observer specialised in maintaining
  distinctions under repeated, intersubjectively verified perturbation.
  Scientific progress is the expansion of $\Rest(\mathcal{R}_{\mathrm{lab}})$:
  the set of distinctions that the collective repair apparatus of
  science can maintain.
  A scientific \emph{revolution} (in Kuhn's sense) is a restructuring
  of $\mathcal{R}_{\mathrm{lab}}$ that changes which distinctions are
  restorable, often making previously central distinctions
  unrestorable while opening new regions of restoration.
  %% EXPAND: connection to Kuhn, paradigm shift as repair-repertoire change.
\end{example}

\begin{theorem}[Paradigm Shift as Repair-Repertoire Discontinuity]
  Let $\mathcal{R}_{\mathrm{old}}$ and $\mathcal{R}_{\mathrm{new}}$
  be two repair repertoires representing a scientific community before
  and after a paradigm shift.
  If the shift is \emph{non-trivial}, meaning
  $\mathcal{R}_{\mathrm{old}} \not\subseteq \mathcal{R}_{\mathrm{new}}$
  and
  $\mathcal{R}_{\mathrm{new}} \not\subseteq \mathcal{R}_{\mathrm{old}}$,
  then the restorable distinction sets are incomparable:
  \[
    \Rest(\mathcal{R}_{\mathrm{old}})
    \;\not\subseteq\;
    \Rest(\mathcal{R}_{\mathrm{new}})
    \qquad\text{and}\qquad
    \Rest(\mathcal{R}_{\mathrm{new}})
    \;\not\subseteq\;
    \Rest(\mathcal{R}_{\mathrm{old}}).
  \]
  In particular, the symmetric difference
  \[
    \Rest(\mathcal{R}_{\mathrm{old}})
    \;\triangle\;
    \Rest(\mathcal{R}_{\mathrm{new}})
    \;\neq\; \emptyset
  \]
  contains \emph{lost distinctions}
  $\Rest(\mathcal{R}_{\mathrm{old}}) \setminus
  \Rest(\mathcal{R}_{\mathrm{new}})$
  and \emph{gained distinctions}
  $\Rest(\mathcal{R}_{\mathrm{new}}) \setminus
  \Rest(\mathcal{R}_{\mathrm{old}})$.
  \label{thm:paradigm}
\end{theorem}

\begin{proof}
  By contraposition of Theorem~\ref{thm:monotonicity}: if
  $\mathcal{R}_{\mathrm{old}} \not\subseteq \mathcal{R}_{\mathrm{new}}$,
  there exists $r \in \mathcal{R}_{\mathrm{old}} \setminus
  \mathcal{R}_{\mathrm{new}}$.
  Any distinction $d$ whose only repair paths under admissible
  perturbations pass through $r$ is in
  $\Rest(\mathcal{R}_{\mathrm{old}}) \setminus
  \Rest(\mathcal{R}_{\mathrm{new}})$; such distinctions exist
  generically when $r$ is not compositionally redundant.
  The symmetric argument yields gains.
\end{proof}

\begin{remark}
  Theorem~\ref{thm:paradigm} formalises Kuhnian incommensurability~\cite{Kuhn1962}
  as a repair-space statement.
  Lost distinctions are not \emph{refuted}; they become
  \emph{unrestorable} in the new paradigm.
  This is precisely the claim that different paradigms make different
  kinds of mistakes visible and invisible, not merely different claims
  about the same visible space.
  The interference matrix $\IntMat$ of Section~\ref{sec:structure}
  will in general have a different spectrum under
  $\mathcal{R}_{\mathrm{old}}$ and $\mathcal{R}_{\mathrm{new}}$,
  providing a quantitative signature of the shift.
  This connects the theorem to the distinction-reorganization papers
  in the broader program.
\end{remark}

\begin{definition}[Social reality as institutional restorability]
  An institutional distinction $d$ is \emph{socially real} relative
  to institutional observer $O_{\mathrm{inst}}$ if and only if
  $d \in \Rest(\mathcal{R}_{\mathrm{inst}})$.
  \label{def:social-reality}
\end{definition}

\begin{proposition}[Institutional collapse shrinks the real]
  Let $O_{\mathrm{inst}}$ undergo a perturbation that reduces its
  repair repertoire from $\mathcal{R}$ to $\mathcal{R}' \subsetneq
  \mathcal{R}$.
  Then
  \[
    \Rest(\mathcal{R}') \;\subseteq\; \Rest(\mathcal{R}),
  \]
  and the set of socially real distinctions is (weakly) smaller after
  the collapse.
  Distinctions in $\Rest(\mathcal{R}) \setminus \Rest(\mathcal{R}')$
  do not become false; they become unrestorable.
  \label{prop:collapse}
\end{proposition}

\begin{proof}
  By Theorem~\ref{thm:monotonicity} applied contravariantly:
  $\mathcal{R}' \subseteq \mathcal{R}$ implies
  $\Rest(\mathcal{R}') \subseteq \Rest(\mathcal{R})$.
  Any distinction in the difference set was formerly restorable
  by procedures no longer available; it is now unrestorable, hence
  socially unreal by Definition~\ref{def:social-reality}.
\end{proof}


% ============================================================
\section{Consequences and Open Problems}
\label{sec:conclusion}
% ============================================================

Taking $\Obs(O) = \Rest(O)$ seriously reorganises several independent
lines of inquiry into a single structure.

\medskip\noindent
\textbf{Memory.}\;
Memory is not storage of representations but maintenance of repair
trajectories.
To remember a distinction is to retain a repair path by which it can
be restored after perturbation.
Forgetting is the loss of repair paths, not the erasure of stored
data.
This connects to the deformation-of-the-admissibility-manifold account
of memory loss developed separately in the \emph{Cost of Forgetting}
program.

\medskip\noindent
\textbf{Understanding.}\;
Understanding a distinction is possessing an efficient repair path
for it.
Understanding increases with the depth and economy of available repair
trajectories, not with the volume of information received.

\medskip\noindent
\textbf{Intelligence.}\;
Intelligence, under this framework, is repair efficiency: the capacity
to restore distinctions with minimal cost across the widest possible
range of perturbations.
The admissibility manifold of an intelligent system is larger and more
robustly bounded than that of a less intelligent system.

\medskip\noindent
\textbf{Science.}\;
Science is collective repair under intersubjective constraint.
The objectivity of scientific knowledge consists not in its
correspondence to an observer-independent world but in its
intersubjective restorability: the distinctions that science maintains
are those restorable by the collective repair apparatus, reproducibly,
across observers.

\medskip
The following problems remain open.

\begin{enumerate}
  \item \textbf{Limit restorability theorem.}\;
    Theorem~\ref{thm:stable-obs} characterises $\Restf(\mathcal{R})$
    via stable observability.
    Prove the analogue for $\Reston(\mathcal{R})$ under weaker
    topological assumptions: does there exist a pre-theoretic notion
    of \emph{asymptotic observability} that identifies exactly with
    limit restorability?

  \item \textbf{Interference matrix spectrum.}\;
    Characterise which symmetric real matrices arise as interference
    matrices $\IntMat$ of some repair repertoire.
    What constraints does the repair structure impose on the spectrum?
    Is there a repair analogue of the Wigner semicircle law for
    large randomly generated observers?

  \item \textbf{Repair conflict graph.}\;
    The positive part of $\IntMat$ defines a signed graph on
    $\Rest(\mathcal{R})$ whose positive-weight cliques are repair
    bottleneck clusters.
    Characterise graphs arising from repair repertoires.
    Relate the clique number to attentional capacity limits in
    biological observers.
    The empirical finding that production embedding models concentrate
    variance into roughly 16 effective dimensions regardless of nominal
    dimensionality~\cite{RayBarman2026} suggests that biological and
    artificial retrieval systems may share a characteristic interference
    graph structure; characterising this structure in repair-theoretic
    terms is an open problem.

  \item \textbf{Repair basis.}\;
    Is there a notion of \emph{repair basis} for a distinction space,
    analogous to a basis for a vector space?
    A natural candidate: a minimal set $B \subseteq \Rest(\mathcal{R})$
    such that every element of $\Rest(\mathcal{R})$ is limit-restorable
    from synergistic combinations of elements of $B$.

  \item \textbf{Repair dynamics and manifold deformation.}\;
    How does $\AO$ evolve when $\mathcal{R}_O$ is itself perturbed?
    Is manifold deformation under memory loss (the Cost of Forgetting
    program) equivalent to repair-repertoire reduction
    $\mathcal{R} \to \mathcal{R}' \subsetneq \mathcal{R}$?
    Is the resulting deformation of $\Rest(\mathcal{R})$ under $d_R$
    an isometric embedding or a metric collapse?

  \item \textbf{Control--Geometry Collapse without compactness.}\;
    Theorem~\ref{thm:cg-collapse} assumes $\Xsub$ compact and convex.
    What are the weakest conditions under which
    $\Rest(\mathcal{R}_O) = \Reach(\AO)$?
    Characterise the Lie algebra rank conditions on $\mathcal{R}_O$
    under which the inclusion collapses, and map this onto
    topological controllability of RSVP field profiles.

  \item \textbf{Interface theorems.}\;
    Prove a precise connection between the Repair Boundary Theorem
    (Theorem~\ref{thm:repair-boundary}) and the interface-derivation
    program in quadratic gravity.
    Under the duality of Section~\ref{sec:structure}, the
    quantum-classical interface is the locus where quantum repair
    operators cease to be available at the classical substrate, i.e.,
    where $\CR$ diverges for quantum state-preparation distinctions.
    Make this precise using the Reachability--Restorability framework.
\end{enumerate}

\medskip
Observation is not passive reception.
It is the maintenance of a world under perturbation.
The world an observer inhabits is exactly as large as its capacity to
repair.

% ============================================================
\begin{thebibliography}{99}

\bibitem{Ashby1956}
Ashby, W.~Ross.
\newblock \emph{An Introduction to Cybernetics}.
\newblock Chapman \& Hall, London, 1956.

\bibitem{ConantAshby1970}
Conant, Roger~C., and W.~Ross Ashby.
\newblock Every good regulator of a system must be a model of that system.
\newblock \emph{International Journal of Systems Science}, 1(2):89--97, 1970.

\bibitem{MaturanaVarela1980}
Maturana, Humberto~R., and Francisco~J. Varela.
\newblock \emph{Autopoiesis and Cognition: The Realization of the Living}.
\newblock D.~Reidel, Dordrecht, 1980.

\bibitem{VarelaThompsonRosch1991}
Varela, Francisco~J., Evan Thompson, and Eleanor Rosch.
\newblock \emph{The Embodied Mind}.
\newblock MIT Press, Cambridge, MA, 1991.

\bibitem{Gibson1979}
Gibson, James~J.
\newblock \emph{The Ecological Approach to Visual Perception}.
\newblock Houghton Mifflin, Boston, 1979.

\bibitem{Marr1982}
Marr, David.
\newblock \emph{Vision}.
\newblock W.~H. Freeman, San Francisco, 1982.

\bibitem{Weiskrantz1986}
Weiskrantz, Lawrence.
\newblock \emph{Blindsight: A Case Study and Implications}.
\newblock Oxford University Press, Oxford, 1986.

\bibitem{MilnerGoodale2006}
Milner, A.~David, and Melvyn~A. Goodale.
\newblock \emph{The Visual Brain in Action}.
\newblock Oxford University Press, Oxford, 2nd edition, 2006.

\bibitem{Kuhn1962}
Kuhn, Thomas~S.
\newblock \emph{The Structure of Scientific Revolutions}.
\newblock University of Chicago Press, Chicago, 1962.

\bibitem{Bellman1957}
Bellman, Richard.
\newblock \emph{Dynamic Programming}.
\newblock Princeton University Press, Princeton, NJ, 1957.

\bibitem{Kalman1960}
Kalman, Rudolf~E.
\newblock Contributions to the theory of optimal control.
\newblock \emph{Boletin de la Sociedad Matem\'atica Mexicana},
  5:102--119, 1960.

\bibitem{Sontag1998}
Sontag, Eduardo~D.
\newblock \emph{Mathematical Control Theory}.
\newblock Springer, New York, 1998.

\bibitem{Munkres2000}
Munkres, James~R.
\newblock \emph{Topology}.
\newblock Prentice Hall, Upper Saddle River, NJ, 2nd edition, 2000.

\bibitem{RockafellarWets1998}
Rockafellar, R.~Tyrrell, and Roger~J.-B. Wets.
\newblock \emph{Variational Analysis}.
\newblock Springer, Berlin, 1998.

\bibitem{Villani2009}
Villani, C\'edric.
\newblock \emph{Optimal Transport: Old and New}.
\newblock Springer, Berlin, 2009.

\bibitem{Shannon1948}
Shannon, Claude~E.
\newblock A mathematical theory of communication.
\newblock \emph{Bell System Technical Journal}, 27(3):379--423, 1948.

\bibitem{Rovelli1996}
Rovelli, Carlo.
\newblock Relational quantum mechanics.
\newblock \emph{International Journal of Theoretical Physics},
  35(8):1637--1678, 1996.

\bibitem{Zurek2003}
Zurek, Wojciech~H.
\newblock Decoherence, einselection, and the quantum origins of the classical.
\newblock \emph{Reviews of Modern Physics}, 75(3):715--775, 2003.

\bibitem{RayBarman2026}
Ray~Barman, Sambartha, Andrey Starenky, Sophia Bodnar, Nikhil Narasimhan,
  and Ashwin Gopinath.
\newblock The geometry of forgetting.
\newblock \emph{arXiv}:2604.06222, 2026.

\end{thebibliography}

\end{document}
